\chapter*{ضمیمهٔ ۲: شخصیت‌نامه}
\addcontentsline{toc}{chapter}{ضمیمهٔ ۲: شخصیت‌نامه}

\begin{center}
\textit{صد چهرهٔ کلیدی در تاریخ اندیشهٔ سیاسی}
\end{center}

\vspace{5mm}

\section*{فیلسوفان و نظریه‌پردازان}

\begin{longtable}{|>{\columncolor{leftred!10}}p{2.5cm}|p{1.5cm}|p{2cm}|p{3cm}|p{5cm}|}
\hline
\rowcolor{leftred!30}
\textbf{نام} & \textbf{سال‌ها} & \textbf{ملیت} & \textbf{گرایش} & \textbf{اهمیت} \\
\hline
\endfirsthead
\hline
\rowcolor{leftred!30}
\textbf{نام} & \textbf{سال‌ها} & \textbf{ملیت} & \textbf{گرایش} & \textbf{اهمیت} \\
\hline
\endhead

کارل مارکس & ۱۸۱۸-۱۸۸۳ & آلمانی & سوسیالیسم علمی & بنیان‌گذار مارکسیسم؛ \textit{سرمایه}، \textit{مانیفست} \\
\hline
فردریش انگلس & ۱۸۲۰-۱۸۹۵ & آلمانی & سوسیالیسم علمی & همکار مارکس؛ تدوین‌کنندهٔ مارکسیسم \\
\hline
ادموند برک & ۱۷۲۹-۱۷۹۷ & ایرلندی-بریتانیایی & محافظه‌کاری & پدر محافظه‌کاری مدرن \\
\hline
جان لاک & ۱۶۳۲-۱۷۰۴ & انگلیسی & لیبرالیسم کلاسیک & پدر لیبرالیسم؛ حقوق طبیعی \\
\hline
ژان ژاک روسو & ۱۷۱۲-۱۷۷۸ & سوئیسی-فرانسوی & جمهوری‌خواهی رادیکال & قرارداد اجتماعی؛ ارادهٔ عمومی \\
\hline
آدام اسمیت & ۱۷۲۳-۱۷۹۰ & اسکاتلندی & لیبرالیسم اقتصادی & پدر اقتصاد مدرن؛ \textit{ثروت ملل} \\
\hline
جان استوارت میل & ۱۸۰۶-۱۸۷۳ & انگلیسی & لیبرالیسم & \textit{دربارهٔ آزادی}؛ فایده‌گرایی \\
\hline
فردریش هایک & ۱۸۹۹-۱۹۹۲ & اتریشی-بریتانیایی & لیبرالیسم کلاسیک & \textit{راه بندگی}؛ نوبل اقتصاد ۱۹۷۴ \\
\hline
لودویگ فون میزس & ۱۸۸۱-۱۹۷۳ & اتریشی-آمریکایی & لیبرتارینیسم & مکتب اتریش؛ نقد سوسیالیسم \\
\hline
آین رند & ۱۹۰۵-۱۹۸۲ & روسی-آمریکایی & ابژکتیویسم & \textit{اطلس شانه بالا انداخت} \\
\hline
رابرت نوزیک & ۱۹۳۸-۲۰۰۲ & آمریکایی & لیبرتارینیسم & \textit{آنارشی، دولت و آرمانشهر} \\
\hline
موری راتبارد & ۱۹۲۶-۱۹۹۵ & آمریکایی & آنارکو-کاپیتالیسم & \textit{اخلاق آزادی} \\
\hline
جان رالز & ۱۹۲۱-۲۰۰۲ & آمریکایی & لیبرالیسم برابری‌خواه & \textit{نظریهٔ عدالت} \\
\hline
آنتونیو گرامشی & ۱۸۹۱-۱۹۳۷ & ایتالیایی & مارکسیسم & هژمونی فرهنگی \\
\hline
هانا آرنت & ۱۹۰۶-۱۹۷۵ & آلمانی-آمریکایی & نظریهٔ سیاسی & \textit{توتالیتاریسم}؛ ابتذال شر \\
\hline
میشل فوکو & ۱۹۲۶-۱۹۸۴ & فرانسوی & پساساختارگرایی & قدرت و دانش \\
\hline
آنتونی گیدنز & ۱۹۳۸- & بریتانیایی & راه سوم & مشاور بلر؛ \textit{فراسوی چپ و راست} \\
\hline
مایکل سندل & ۱۹۵۳- & آمریکایی & کمونیتارینیسم & نقد لیبرالیسم؛ فلسفهٔ عمومی \\
\hline
\end{longtable}

\section*{رهبران سیاسی}

\begin{longtable}{|>{\columncolor{rightblue!10}}p{2.5cm}|p{1.5cm}|p{2cm}|p{3cm}|p{5cm}|}
\hline
\rowcolor{rightblue!30}
\textbf{نام} & \textbf{سال‌ها} & \textbf{ملیت} & \textbf{گرایش} & \textbf{اهمیت} \\
\hline
\endfirsthead
\hline
\rowcolor{rightblue!30}
\textbf{نام} & \textbf{سال‌ها} & \textbf{ملیت} & \textbf{گرایش} & \textbf{اهمیت} \\
\hline
\endhead

ولادیمیر لنین & ۱۸۷۰-۱۹۲۴ & روسی & کمونیسم & رهبر انقلاب روسیه؛ بنیان‌گذار شوروی \\
\hline
لئون تروتسکی & ۱۸۷۹-۱۹۴۰ & روسی & کمونیسم & انقلاب مداوم؛ رقیب استالین \\
\hline
ژوزف استالین & ۱۸۷۸-۱۹۵۳ & گرجی-روسی & کمونیسم & رهبر شوروی؛ سوسیالیسم در یک کشور \\
\hline
مائو تسه‌تونگ & ۱۸۹۳-۱۹۷۶ & چینی & کمونیسم & رهبر چین کمونیست؛ مائوئیسم \\
\hline
بنیتو موسولینی & ۱۸۸۳-۱۹۴۵ & ایتالیایی & فاشیسم & بنیان‌گذار فاشیسم \\
\hline
آدولف هیتلر & ۱۸۸۹-۱۹۴۵ & اتریشی-آلمانی & نازیسم & رهبر آلمان نازی؛ هولوکاست \\
\hline
وینستون چرچیل & ۱۸۷۴-۱۹۶۵ & بریتانیایی & محافظه‌کاری & نخست‌وزیر جنگ؛ نماد مقاومت \\
\hline
فرانکلین روزولت & ۱۸۸۲-۱۹۴۵ & آمریکایی & لیبرالیسم جدید & نیو دیل؛ چهار دوره ریاست‌جمهوری \\
\hline
شارل دوگل & ۱۸۹۰-۱۹۷۰ & فرانسوی & گلیسم & جمهوری پنجم؛ استقلال‌طلبی \\
\hline
ویلی برانت & ۱۹۱۳-۱۹۹۲ & آلمانی & سوسیال دموکراسی & صدراعظم آلمان؛ سیاست شرق \\
\hline
مارگارت تاچر & ۱۹۲۵-۲۰۱۳ & بریتانیایی & نئولیبرالیسم & انقلاب تاچری؛ «بانوی آهنین» \\
\hline
رونالد ریگان & ۱۹۱۱-۲۰۰۴ & آمریکایی & نئولیبرالیسم & انقلاب ریگانی؛ پایان جنگ سرد \\
\hline
میخائیل گورباچف & ۱۹۳۱-۲۰۲۲ & روسی & اصلاح‌طلبی & گلاسنوست و پرسترویکا \\
\hline
نلسون ماندلا & ۱۹۱۸-۲۰۱۳ & آفریقای جنوبی & سوسیال دموکراسی & پایان آپارتاید؛ آشتی ملی \\
\hline
تونی بلر & ۱۹۵۳- & بریتانیایی & راه سوم & نیو لیبر؛ سه دورهٔ نخست‌وزیری \\
\hline
باراک اوباما & ۱۹۶۱- & آمریکایی & لیبرالیسم & اولین رئیس‌جمهور آفریقایی‌تبار \\
\hline
دونالد ترامپ & ۱۹۴۶- & آمریکایی & پوپولیسم راست & MAGA؛ ناسیونالیسم اقتصادی \\
\hline
ویکتور اوربان & ۱۹۶۳- & مجارستانی & پوپولیسم راست & «دموکراسی غیرلیبرال» \\
\hline
\end{longtable}

\section*{اقتصاددانان}

\begin{longtable}{|>{\columncolor{goldaccent!10}}p{2.5cm}|p{1.5cm}|p{2cm}|p{3cm}|p{5cm}|}
\hline
\rowcolor{goldaccent!30}
\textbf{نام} & \textbf{سال‌ها} & \textbf{ملیت} & \textbf{گرایش} & \textbf{اهمیت} \\
\hline
\endfirsthead

جان مینارد کینز & ۱۸۸۳-۱۹۴۶ & بریتانیایی & کینزگرایی & انقلاب کینزی؛ دولت رفاه \\
\hline
میلتون فریدمن & ۱۹۱۲-۲۰۰۶ & آمریکایی & مونتاریسم & مکتب شیکاگو؛ نوبل ۱۹۷۶ \\
\hline
جوزف استیگلیتز & ۱۹۴۳- & آمریکایی & کینزگرایی جدید & نقد جهانی‌شدن؛ نوبل ۲۰۰۱ \\
\hline
پل کروگمن & ۱۹۵۳- & آمریکایی & لیبرالیسم & نیویورک تایمز؛ نوبل ۲۰۰۸ \\
\hline
توماس پیکتی & ۱۹۷۱- & فرانسوی & چپ & \textit{سرمایه در قرن بیست‌ویکم} \\
\hline
\end{longtable}

\section*{فعالان و روشنفکران}

\begin{longtable}{|>{\columncolor{centergreen!10}}p{2.5cm}|p{1.5cm}|p{2cm}|p{3cm}|p{5cm}|}
\hline
\rowcolor{centergreen!30}
\textbf{نام} & \textbf{سال‌ها} & \textbf{ملیت} & \textbf{گرایش} & \textbf{اهمیت} \\
\hline
\endfirsthead

رزا لوکزامبورگ & ۱۸۷۱-۱۹۱۹ & لهستانی-آلمانی & سوسیالیسم انقلابی & نقد لنینیسم و رفرمیسم \\
\hline
ارنستو چه گوارا & ۱۹۲۸-۱۹۶۷ & آرژانتینی-کوبایی & کمونیسم & انقلاب کوبا؛ نماد چپ \\
\hline
مارتین لوتر کینگ & ۱۹۲۹-۱۹۶۸ & آمریکایی & حقوق مدنی & جنبش حقوق مدنی؛ ضد خشونت \\
\hline
سیمون دوبووار & ۱۹۰۸-۱۹۸۶ & فرانسوی & فمینیسم & \textit{جنس دوم} \\
\hline
نوام چامسکی & ۱۹۲۸- & آمریکایی & آنارکو-سندیکالیسم & زبان‌شناسی؛ نقد سیاست آمریکا \\
\hline
جردن پیترسون & ۱۹۶۲- & کانادایی & محافظه‌کاری & نقد چپ فرهنگی؛ روانشناسی \\
\hline
\end{longtable}

\begin{longtable}{|>{\columncolor{rightblue!10}}p{2.5cm}|p{1.5cm}|p{2cm}|p{3cm}|p{5cm}|}
\hline
\rowcolor{rightblue!30}
\textbf{نام} & \textbf{سال‌ها} & \textbf{ملیت} & \textbf{گرایش} & \textbf{اهمیت} \\
\hline

فیدل کاسترو & ۱۹۲۶-۲۰۱۶ & کوبایی & کمونیسم & انقلاب کوبا؛ ۵۰ سال حکومت \\
\hline
سالوادور آلنده & ۱۹۰۸-۱۹۷۳ & شیلیایی & سوسیالیسم & اولین رئیس‌جمهور سوسیالیست منتخب \\
\hline
اوگوستو پینوشه & ۱۹۱۵-۲۰۰۶ & شیلیایی & اقتدارگرایی راست & کودتا ۱۹۷۳؛ آزمایشگاه نئولیبرالیسم \\
\hline
جمال عبدالناصر & ۱۹۱۸-۱۹۷۰ & مصری & ناسیونالیسم عربی & ناصریسم؛ ملی‌سازی سوئز \\
\hline
آیت‌الله خمینی & ۱۹۰۲-۱۹۸۹ & ایرانی & اسلام‌گرایی & انقلاب اسلامی؛ ولایت فقیه \\
\hline
محمد مصدق & ۱۸۸۲-۱۹۶۷ & ایرانی & ناسیونالیسم لیبرال & ملی‌سازی نفت؛ کودتای ۱۳۳۲ \\
\hline
جواهرلعل نهرو & ۱۸۸۹-۱۹۶۴ & هندی & سوسیالیسم دموکراتیک & بنیان‌گذار هند مدرن؛ عدم تعهد \\
\hline
نارندرا مودی & ۱۹۵۰- & هندی & ناسیونالیسم هندو & نخست‌وزیر از ۲۰۱۴؛ BJP \\
\hline
مائو تسه‌تونگ & ۱۸۹۳-۱۹۷۶ & چینی & کمونیسم & انقلاب چین؛ انقلاب فرهنگی \\
\hline
دنگ شیائوپینگ & ۱۹۰۴-۱۹۹۷ & چینی & سوسیالیسم بازار & اصلاح و گشایش؛ معجزهٔ اقتصادی \\
\hline
شی جین‌پینگ & ۱۹۵۳- & چینی & اقتدارگرایی & رهبر چین از ۲۰۱۲؛ قدرت مادام‌العمر \\
\hline
لی کوان یو & ۱۹۲۳-۲۰۱۵ & سنگاپوری & اقتدارگرایی توسعه‌گرا & معمار سنگاپور مدرن \\
\hline
لولا داسیلوا & ۱۹۴۵- & برزیلی & چپ میانه & رئیس‌جمهور دو دوره؛ موج صورتی \\
\hline
ژائیر بولسونارو & ۱۹۵۵- & برزیلی & راست پوپولیست & «ترامپ برزیل»؛ رئیس‌جمهور ۲۰۱۹-۲۰۲۲ \\
\hline
خاویر میلی & ۱۹۷۰- & آرژانتینی & لیبرتارینیسم افراطی & رئیس‌جمهور از ۲۰۲۳ \\
\hline
فرانسوا میتران & ۱۹۱۶-۱۹۹۶ & فرانسوی & سوسیالیسم & رئیس‌جمهور ۱۹۸۱-۱۹۹۵ \\
\hline
امانوئل ماکرون & ۱۹۷۷- & فرانسوی & لیبرالیسم میانه & «نه چپ، نه راست» \\
\hline
آنگلا مرکل & ۱۹۵۴- & آلمانی & دموکراسی مسیحی & صدراعظم ۲۰۰۵-۲۰۲۱ \\
\hline
اولاف شولتس & ۱۹۵۸- & آلمانی & سوسیال دموکراسی & صدراعظم از ۲۰۲۱ \\
\hline
سیلویو برلوسکونی & ۱۹۳۶-۲۰۲۳ & ایتالیایی & پوپولیسم راست & نخست‌وزیر سه دوره \\
\hline
جورجا ملونی & ۱۹۷۷- & ایتالیایی & راست رادیکال & نخست‌وزیر از ۲۰۲۲ \\
\hline
\end{longtable}

\section*{متفکران و فعالان قرن بیستم و بیست‌ویکم}

\begin{longtable}{|>{\columncolor{violet!10}}p{2.5cm}|p{1.5cm}|p{2cm}|p{3cm}|p{5cm}|}
\hline
\rowcolor{violet!30}
\textbf{نام} & \textbf{سال‌ها} & \textbf{ملیت} & \textbf{گرایش} & \textbf{اهمیت} \\
\hline

هربرت مارکوزه & ۱۸۹۸-۱۹۷۹ & آلمانی-آمریکایی & نئومارکسیسم & مکتب فرانکفورت؛ چپ نو \\
\hline
فرانتس فانون & ۱۹۲۵-۱۹۶۱ & مارتینیکی & ضد استعمار & \textit{دوزخیان زمین}؛ الجزایر \\
\hline
پائولو فریره & ۱۹۲۱-۱۹۹۷ & برزیلی & آموزش انتقادی & \textit{آموزش ستمدیدگان} \\
\hline
سیمون دوبووار & ۱۹۰۸-۱۹۸۶ & فرانسوی & فمینیسم & \textit{جنس دوم} \\
\hline
بتی فریدان & ۱۹۲۱-۲۰۰۶ & آمریکایی & فمینیسم & \textit{رمز و راز زنانه}؛ NOW \\
\hline
جودیت باتلر & ۱۹۵۶- & آمریکایی & پساساختارگرایی فمینیستی & نظریهٔ جنسیت \\
\hline
کارل اشمیت & ۱۸۸۸-۱۹۸۵ & آلمانی & محافظه‌کاری اقتدارگرا & نظریهٔ دولت؛ نقد لیبرالیسم \\
\hline
لئو اشتراوس & ۱۸۹۹-۱۹۷۳ & آلمانی-آمریکایی & محافظه‌کاری & فلسفهٔ سیاسی؛ نئوکنسرواتیسم \\
\hline
آلن بلوم & ۱۹۳۰-۱۹۹۲ & آمریکایی & محافظه‌کاری & \textit{بسته شدن ذهن آمریکایی} \\
\hline
راجر اسکروتن & ۱۹۴۴-۲۰۲۰ & بریتانیایی & محافظه‌کاری & فیلسوف محافظه‌کار معاصر \\
\hline
مورای بوکچین & ۱۹۲۱-۲۰۰۶ & آمریکایی & شهرداری آزادی‌خواه & اکولوژی اجتماعی \\
\hline
دیوید گرایبر & ۱۹۶۱-۲۰۲۰ & آمریکایی & آنارشیسم & جنبش اشغال وال‌استریت \\
\hline
نائومی کلاین & ۱۹۷۰- & کانادایی & چپ & \textit{دکترین شوک}؛ اقلیم \\
\hline
توماس سوول & ۱۹۳۰- & آمریکایی & لیبرتارینیسم & اقتصاددان محافظه‌کار \\
\hline
ویلیام باکلی & ۱۹۲۵-۲۰۰۸ & آمریکایی & محافظه‌کاری & National Review؛ جنبش محافظه‌کار \\
\hline
استیو بنن & ۱۹۵۳- & آمریکایی & ناسیونال-پوپولیسم & برایتبارت؛ مشاور ترامپ \\
\hline
ریچارد اسپنسر & ۱۹۷۸- & آمریکایی & آلت‌رایت & ناسیونالیسم سفید \\
\hline
جردن پیترسون & ۱۹۶۲- & کانادایی & محافظه‌کاری & نقد چپ فرهنگی؛ روانشناسی \\
\hline
بن شاپیرو & ۱۹۸۴- & آمریکایی & محافظه‌کاری & رسانهٔ محافظه‌کار جوان \\
\hline
سلاوی ژیژک & ۱۹۴۹- & اسلوونیایی & مارکسیسم & فیلسوف فرهنگی؛ نقد لیبرالیسم \\
\hline
یورگن هابرماس & ۱۹۲۹- & آلمانی & نظریهٔ انتقادی & کنش ارتباطی؛ حوزهٔ عمومی \\
\hline
\end{longtable}

\section*{شخصیت‌های تاریخی بنیادی}

\begin{longtable}{|>{\columncolor{goldaccent!10}}p{2.5cm}|p{1.5cm}|p{2cm}|p{3cm}|p{5cm}|}
\hline
\rowcolor{goldaccent!30}
\textbf{نام} & \textbf{سال‌ها} & \textbf{ملیت} & \textbf{گرایش} & \textbf{اهمیت} \\
\hline

ماکیاولی & ۱۴۶۹-۱۵۲۷ & ایتالیایی & رئالیسم سیاسی & \textit{شهریار}؛ پدر علم سیاست مدرن \\
\hline
توماس مور & ۱۴۷۸-۱۵۳۵ & انگلیسی & اومانیسم & \textit{آرمانشهر} \\
\hline
توماس هابز & ۱۵۸۸-۱۶۷۹ & انگلیسی & قرارداد اجتماعی & \textit{لویاتان}؛ حاکمیت مطلق \\
\hline
جان لاک & ۱۶۳۲-۱۷۰۴ & انگلیسی & لیبرالیسم & حقوق طبیعی؛ حکومت مشروطه \\
\hline
مونتسکیو & ۱۶۸۹-۱۷۵۵ & فرانسوی & لیبرالیسم & تفکیک قوا؛ \textit{روح القوانین} \\
\hline
ژان ژاک روسو & ۱۷۱۲-۱۷۷۸ & سوئیسی-فرانسوی & جمهوری‌خواهی & قرارداد اجتماعی؛ ارادهٔ عمومی \\
\hline
آدام اسمیت & ۱۷۲۳-۱۷۹۰ & اسکاتلندی & لیبرالیسم اقتصادی & \textit{ثروت ملل}؛ دست نامرئی \\
\hline
ادموند برک & ۱۷۲۹-۱۷۹۷ & ایرلندی & محافظه‌کاری & نقد انقلاب فرانسه \\
\hline
توماس پین & ۱۷۳۷-۱۸۰۹ & انگلیسی-آمریکایی & جمهوری‌خواهی رادیکال & \textit{حقوق بشر}؛ \textit{عقل سلیم} \\
\hline
ماری ولستونکرافت & ۱۷۵۹-۱۷۹۷ & انگلیسی & فمینیسم اولیه & \textit{احقاق حقوق زن} \\
\hline
جرمی بنتام & ۱۷۴۸-۱۸۳۲ & انگلیسی & فایده‌گرایی & بیشترین خوشبختی برای بیشترین تعداد \\
\hline
الکسی دو توکویل & ۱۸۰۵-۱۸۵۹ & فرانسوی & لیبرالیسم محافظه‌کارانه & \textit{دموکراسی در آمریکا} \\
\hline
جوزپه ماتزینی & ۱۸۰۵-۱۸۷۲ & ایتالیایی & ناسیونالیسم جمهوری‌خواه & وحدت ایتالیا \\
\hline
هنری دیوید ثورو & ۱۸۱۷-۱۸۶۲ & آمریکایی & فردگرایی آنارشیستی & نافرمانی مدنی \\
\hline
فردریش نیچه & ۱۸۴۴-۱۹۰۰ & آلمانی & فراسوی خیر و شر & نقد اخلاق؛ ابرانسان \\
\hline
\end{longtable}